\chapter{长期规划：滚动承诺、模型修订与目标治理}
\label{chp:long_term_planning_and_goal_revision}
长期规划不是由某个未来状态逆向施加因果力，也不是让智能体在一次计算中枚举所有可能路线。它要求系统把远期规范转化为当前可执行、可验证且可撤回的阶段承诺，并在环境证据改变时区分三件事：调整动作、修订世界模型、还是依法改变目标。本章承接目标—身份—执行控制的分工，以部分可观测的混合状态、滚动时域优化、版本化世界模型和目标治理描述这一过程。我们保留去罗马、断桥、绕路与转移目的地等例证，但把“未来吸引子”“地形重绘”和“奇点蒸发”重建为明确的计算接口，使计划成功、内部代价下降、资源消耗与真实环境结果各自可检验。

\section{长期规划的对象：从未来约束到滚动时域承诺}
\label{sec:physics_of_planning}
我们先定义规划处理的对象。离散决策时刻为 $k\in\mathbb N$，环境真实状态 $x_k\in\mathcal X$ 通常不可完全观测，系统维护信念 $b_k\in\mathcal P(\mathcal X)$；动作 $u_k\in\mathcal U$ 通过环境转移核 $P(x_{k+1}\mid x_k,u_k)$ 产生后果，观测 $o_{k+1}\in\mathcal O$ 再经似然模型更新信念。任务规范 $\Theta_k=(\mathcal C_k,J_k,\mathcal R_k,\nu_k)$ 分别记录硬约束、阶段与终端代价、风险度量及版本来源。计划不是一条脱离这些对象的几何曲线，而是带条件分支、资源预算和验证节点的策略片段 $\Pi_{k:k+H}$。状态、信念、动作、代价和实际时间采用不同单位，必须先校准再进入同一优化问题。

未来目标对当前选择的影响是规范性的而非逆因果的。截止日期、交付质量或目的地由当前记录中的目标版本表示；求解器在时刻 $k$ 读取它们，比较不同动作对未来结果的预测。所谓“未来牵引现在”，只说明终端条件参与当前评分，不说明未来事件向过去传播信号。若目标尚未被授权、终端观测不可辨识或预测模型没有覆盖关键变量，系统最多生成条件计划，不能把想象的终点当作已存在事实。这个区分保留了边值问题的数学价值，也避免用块宇宙或超前势替代数据流和因果顺序。

在长度为 $H$ 的预测窗口内，一个风险敏感滚动规划器可求解
\[
\min_{\pi_{k:k+H-1}}\;\mathbb E\!\left[\sum_{j=0}^{H-1}
\ell(b_{k+j},u_{k+j};\Theta_{k+j})+\phi(b_{k+H};\Theta_{k+H})\right]
+\lambda\operatorname{Risk}(\Pi),
\]
同时满足动作权限、资源、时限和安全约束。$\ell$ 是阶段任务代价，$\phi$ 是终端条件的软惩罚，$\operatorname{Risk}$ 只在声明的概率模型和尾部事件定义下成立。若目标、输入、模型或度量随时间改变，它们都要带索引进入优化；若使用连续近似，还需声明积分器与离散步长。该问题的最优解只是在当前模型和窗口内的候选，不是全局最短路径，更不证明任务必然成功。

系统只提交计划的第一步或一个短的可撤销片段，收到环境回执后再更新信念、缩短或平移窗口并重新求解。这样既避免一次性展开完整搜索树，也不声称“无需搜索”。候选生成可以采用图搜索、采样、梯度优化、启发式分解或模型预测控制；流式分布表示适合生成相似路线，局部结构网络适合保存地点身份、前置关系、权限和不可违反的依赖。二者通过同一任务键绑定。任何算法都需付出搜索、模拟和验证成本，所谓路径自然涌现只是某些求解器的直观描述。

以去罗马为例，当前信念包含位置、交通状态和信息时效，目标版本记录到达时间、预算和允许的交通方式。规划器可产生经桥梁、铁路或航班的多条条件路线，并为每条路线标注票据、换乘、风险和需要确认的外部事件。它不会把“罗马”设成负无穷势阱，因为再高的软优先级也不能覆盖安全、合法性和资源硬约束；也不会在一次提交中锁死所有后续动作。若第一段道路拥堵，系统可以只改战术动作；若交通模型系统性失准，则进入模型修订；若目的本身失去授权或代价超出容许边界，才进入目标治理。

长期性来自跨窗口保持约束和承诺，而不来自计划无限延伸。我们维护承诺集 $K_k$，每项含责任主体、适用域、截止条件、证据和退出条款。滚动规划不能为了降低当前代价而悄悄丢弃未来承诺；若确需冲突仲裁，必须给出被牺牲目标、影响范围和批准者。反过来，过期承诺也不能永久占据资源。计划缓存因此按目标版本、模型版本和身份键索引，任一关键版本改变都触发失效检查，而不是无条件复用旧路线。

验收包含四类反例。第一，终点相同但身份权限不同，路线应随可用动作集改变；第二，内部代价相同但环境风险不同，系统应保留风险差异；第三，短期代价下降却破坏长期承诺，提交应被阻止或升级审批；第四，模型预测到达而环境没有有效回执，状态不得标记完成。我们同时报告计划可行率、重规划次数、候选覆盖、验证成本、资源消耗和真实成功率。只有在给定模型、扰动和时间窗内这些指标成立，才能说远期规范被可靠地转化为当前承诺。

\section{过程学习：障碍证据、模型版本与受控重规划}
\label{sec:learning_in_process_cefe}
计划执行中的意外首先是观测事件，不是自动改变“认知地形”的物理冲击。令执行记录 $e_k=(u_k,o_{k+1},y_{k+1},\rho_k)$，其中 $y_{k+1}$ 是工具或环境回执，$\rho_k$ 记录来源、时间、完整性和授权。预测残差 $r_k=d(o_{k+1},\widehat o_{k+1})$ 需要一个声明过尺度的比较函数 $d$；残差大只说明模型在该读出上失配，可能来自传感噪声、过期数据、动作未执行或真正障碍。系统不能把一次“惊奇”直接写成永久世界事实，更不能把算法损失、焦虑感受和焦耳能耗合并为一个自由能数值。

断桥例子需要先做身份绑定。系统核验观测指向计划中的哪一座桥、证据是否来自可信交通源、断桥状态的有效时间以及当前路线是否确实依赖该边。如果只收到一条无来源消息，合理动作可以是降低该边置信度、请求第二来源或选择可逆绕行，而不是把距离设为无穷。若现场传感器、官方公告和工具失败日志相互支持，系统生成结构补丁：把对应边的可用状态改为关闭，附上证据、置信度、时间窗和回滚条件。局部结构网络负责桥梁身份与依赖关系，全局分布表示可以提出相似替代路线，但不得用语义相似覆盖对象核验。

模型修订采用候选—验证—提交事务。当前世界模型为 $M_k=(G_k,p_k,\omega_k)$，其中 $G_k$ 是带类型关系图，$p_k$ 是预测组件，$\omega_k$ 是版本和来源元数据。学习器根据事件窗口 $E_{a:b}$ 生成 $\widetilde M=F_M(M_k,E_{a:b})$，随后验证器检查来源一致性、结构约束、预测改进、旧任务回归和安全影响。只有验证报告 $V_M$ 达到门槛，提交器才产生 $M_{k+1}=\operatorname{Commit}(M_k,\widetilde M,V_M)$。图边关闭是离散编辑，预测参数调整是连续更新，两者不能用同一微分方程含混表示。

提交新模型后，旧计划不必全部丢弃。规划器沿依赖图标记受影响的动作和里程碑，只对失效后缀重新求解；未受影响且仍有有效回执的前缀可以保留。增量重规划必须使用幂等动作键，避免工具超时后重复购票、重复付款或重复通知。若所有替代路线都违反约束，求解器返回不可行证书，包括冲突约束、缺失资源和最早失效节点，而不是在内部不断“绕流”。这个证书随后交给目标治理环，不能由模型学习器自行放宽安全边界。

过程学习和规划既耦合又不等同。模型更新改变候选的可达性、转移概率和代价估计，规划结果则决定接下来收集哪些证据；二者形成主动观测闭环。但学习器追求预测改善，规划器追求任务条件下的行动质量，目标函数和失败模式不同。一个更准确的交通模型可能证明罗马在期限内不可达，从而使当前规划代价上升；这仍是学习成功。反之，一次侥幸绕路到达不证明模型正确，因为未识别的错误可能在下一任务暴露。

我们还需控制反馈偏差。若系统只观察自己选择的路线，就缺少对未选路径的证据；若每次障碍都增加永久禁用权重，模型会逐步收缩为保守死区。工程上应保留探索预算、证据过期机制和反事实评估：在不实际执行高风险动作的前提下，用独立数据比较旧模型与候选模型对多个路线的预测。高风险结构更新可先进入影子版本，只有跨窗口复现后才晋级。对真实不可逆后果，内部回滚只能恢复模型，不能修复外部损失，因此验证门槛应随影响增大。

AgentOS实例中，SPC可保存地点、边和证据来源，TCE计算受约束候选，CCM协调模型与目标版本，Boundary阻止越权路径，Offloading与外部记忆 $M_e$ 保存交通回执；$\Psi$ 可作为候选分布，$h(t)$ 表示当前信念摘要。这些名字不是过程学习的普遍必要条件。任何实现只要能维持对象身份、版本、证据、依赖、动作状态和环境验证，就可以实例化同一类计算动力学。

验收通过分层故障注入完成：注入单一坏传感器，系统不应立即永久关闭桥；注入多源一致断桥证据，应只使相关计划后缀失效；制造工具超时，应重试查询而非重复不可逆动作；恢复桥梁并超过有效期，应允许有证据的重新开放；制造真正不可行的全局约束，应把冲突上报目标治理而非擅自改目标。我们报告证据校准、错误结构提交率、受影响范围、重规划延迟、重复动作率和环境成功率，从而把“地形重绘”落实为可审计的模型修订过程。
